\section{The model}

\begin{itemize}
    \item Fixed universal set of existing jobs: $\mathcal{J}$.
    \item Pre-tax income function: $\mathbf{y}:\mathcal{J}\rightarrow \mathbb{R}_{+}$: $\mathbf{y}(j)$, also written $\mathbf{y}_j$, stands for pre-tax income associated to job $j$. 
    \item Consumption set: $Z=\mathbb{R}_{+}\times \mathcal{J}$, typical element $z=(c,j)$, where $c$ is the consumption level and $j$ the job to which the individual is assigned.
    \item Preferences: $R\in\mathcal{R}$, monotonic in the first argument.
    \item Ability set $A\subseteq \mathcal{J}$. The set of ability sets is refer to as $\mathcal{A}$. It can be the entire set of possible subsets of $\mathcal{J}$: $\mathcal{J}:2^{\mathcal{J}}$. We can add restrictions as well, and we introduce these restrictions below when we use them.
    \item Well-being measure: $W: Z\times \mathcal{R} \times 2^{\mathcal{J}}\rightarrow \mathbb{R}$: $W(z,R,A)$ is the well-being level of an individual consuming bundle $z$, having preferences $R$ and ability set $A$.
    \item Well-being is allowed to depend on how much jobs are paid $\mathbf{y}$, so we write $W(z,R,A;\mathbf{y})$.
    \item To save on notation, we assume that as soon as we measure $W(z,R,A;\mathbf{y})$, $z=(c,j)$ for some $j\in A$.
\end{itemize}
Remarks on the model:
\begin{itemize}
  \item By allowing $\mathbf{y}$ to vary, we distinguish between non-pecuniary characteristics of jobs, assumed to be fixed, and pecuniary characteristics, which may vary (across regions or through time). Of course, with a sufficiently large $\mathcal{J}$ we can represent the universe of jobs by fixing $\mathbf{y}$, so that several jobs with exactly the same tasks to perform can be distinguished depending on how much they are paid. We prefer the latter definition, though, to be allowed to make pre-tax income vary without affecting the preferences of individuals towards a particular job.
\end{itemize}
